\documentclass[14pt]{jsarticle}
% ========================================
% 講義MODE プリアンブル（A4 縦 × 2ページ → B4 横に左右2面で縮小印刷）
% 使い方: \documentclass[17pt]{jsarticle}
%         % ========================================
% 講義MODE プリアンブル（A4 縦 × 2ページ → B4 横に左右2面で縮小印刷）
% 使い方: \documentclass[17pt]{jsarticle}
%         % ========================================
% 講義MODE プリアンブル（A4 縦 × 2ページ → B4 横に左右2面で縮小印刷）
% 使い方: \documentclass[17pt]{jsarticle}
%         \input{lecture-preamble.tex}
% ※ B4 に縮小（約87%）すると，17pt は実寸で約14.7pt になる
% ※ jsarticle は \mag で拡大するので，紙面の寸法は truemm で指定する．
%    TikZ の cm も拡大される（1cm → 実寸 約1.7cm）点に注意
% ========================================
\usepackage[dvipdfm, a4paper, margin=15truemm]{geometry}

\input{common-preamble.tex}

\pagestyle{empty}

% ========================================
% ヘッダ（1ページ目の先頭）
% 使い方: \lectureheader{科目}{章番号}{章の名前}{節番号}{節の名前}{項番号}{項の名前}{本時の目標}
% ========================================
\newcommand{\lectureheader}[8]{%
  \setcounter{mathchapter}{#2}%
  \setcounter{section}{#4}%
  \setcounter{subsection}{#6}%
  \noindent\hfill
  年\hspace{2zw}組\hspace{2zw}番\quad 氏名\ \underline{\hspace{8zw}}\par
  \vspace{0.3em}
  \noindent #1\quad 第#2章\quad #3\par
  \noindent{\Large\bfseries 第#4節\quad #5}\hspace{1.5zw}{\large\bfseries #6\quad #7}\par
  \vspace{-0.3em}
  \noindent\rule{\linewidth}{0.8pt}\par
  \vspace{0.3em}
  \begin{itembox}[l]{\textbf{本時の目標}}
    \centering #8
  \end{itembox}
  \vspace{0.3em}
}

% 右のページ（2ページ目）へ移る
\newcommand{\nextpage}{\newpage}

% 見出し: \heading{見出し}
\newcommand{\heading}[1]{%
  \par\noindent{\large\bfseries #1}\par\vspace{0.3em}%
}

% 本時のまとめ
\newenvironment{summarybox}{%
  \begin{itembox}[l]{\textbf{本時のまとめ}}
}{\end{itembox}}

% 左右配置（左：説明・計算，右：図）
% 使い方: \sidebyside[左の幅の割合]{左}{右}
\newcommand{\sidebyside}[3][0.5]{%
  \par\noindent
  \begin{minipage}[t]{#1\linewidth}\vspace{0pt}#2\end{minipage}\hfill
  \begin{minipage}[t]{\dimexpr0.98\linewidth-#1\linewidth\relax}%
    \vspace{0pt}\centering #3\end{minipage}\par
}

% 次回予告（2ページ目の下端）
\newcommand{\nexttime}[1]{%
  \par\noindent\parbox{\linewidth}{%
    \rule{\linewidth}{0.4pt}\par
    次回：#1}\par
}

% ※ B4 に縮小（約87%）すると，17pt は実寸で約14.7pt になる
% ※ jsarticle は \mag で拡大するので，紙面の寸法は truemm で指定する．
%    TikZ の cm も拡大される（1cm → 実寸 約1.7cm）点に注意
% ========================================
\usepackage[dvipdfm, a4paper, margin=15truemm]{geometry}

% ========================================
% 共通プリアンブル（講義MODE・演習MODE）
% パッケージ／解答切り替え／番号／空欄／定義・公式・定理の囲み
% ========================================
\usepackage{amsmath, amssymb, mathtools, bm}
\usepackage{ascmac}
\usepackage[dvipdfmx]{graphicx}
\usepackage[dvipdfmx]{xcolor}
\usepackage{tikz}
\usetikzlibrary{calc, arrows.meta}
\usepackage[most]{tcolorbox}
\usepackage{enumitem}

\linespread{1.25}
\everymath{\displaystyle}
\setlength{\parindent}{0pt}
\setlist[enumerate,1]{label=（\arabic*）, leftmargin=*, labelsep=0.5em}

% ========================================
% 解答表示の切り替え（生徒用：\showanswerfalse ／ 教師用：\showanswertrue）
% ========================================
\newif\ifshowanswer
\showanswertrue

% ========================================
% 番号（「章.節.番号」形式，節ごとにリセット）
% ========================================
% jsarticle には章がないので，章番号は専用カウンタで持つ
\newcounter{mathchapter}
\newcounter{definition}[section]
\newcounter{formula}[section]
\newcounter{theorem}[section]
\renewcommand{\thedefinition}{\arabic{mathchapter}.\arabic{section}.\arabic{definition}}
\renewcommand{\theformula}{\arabic{mathchapter}.\arabic{section}.\arabic{formula}}
\renewcommand{\thetheorem}{\arabic{mathchapter}.\arabic{section}.\arabic{theorem}}

% ========================================
% 書き込み用コマンド
% ========================================
% 空欄（答えの幅＋左右の余白で下線を引く）
% 使い方: \fitblankbf[]{答え}  ※文字列は \text{...} で囲む
\newsavebox{\fitblanksavebox}
\newcommand{\fitblankbf}[2][]{%
  \sbox{\fitblanksavebox}{\ensuremath{\bm{#2}}}%
  \underline{\makebox[\dimexpr\wd\fitblanksavebox+3zw\relax][c]{%
    \ifshowanswer\textcolor{blue}{\usebox{\fitblanksavebox}}%
    \else\phantom{\usebox{\fitblanksavebox}}\fi}}%
}

% 数式中の計算欄（下線なし，答えの幅＋左右2em を確保）
% 使い方: $= \answermath{3 - 1 = 2}$
\newcommand{\answermath}[1]{%
  \ifshowanswer{\color{blue}#1}%
  \else\hspace{2em}\phantom{#1}\hspace{2em}\fi
}

% 記述・計算スペース（薄い枠で高さを確保）
% 使い方: \writespace[高さ]{解答例}
\newcommand{\writespace}[2][5em]{%
  \begin{tcolorbox}[enhanced, colback=white, colframe=black!35,
      boxrule=0.4pt, sharp corners, height=#1, valign=center,
      left=8pt, right=8pt]
    \ifshowanswer\textcolor{blue}{#2}\fi
  \end{tcolorbox}%
}

% ========================================
% 定義（赤系）・公式（青系）・定理（青系・太枠）
% 使い方: \begin{definitionbox}{用語} ... \end{definitionbox}
% ========================================
\tcbset{
  mathbox/.style={
    enhanced, sharp corners,
    left=10pt, right=10pt, top=8pt, bottom=8pt,
    fonttitle=\bfseries,
  }
}
\newtcolorbox{definitionbox}[1]{
  mathbox, colback=red!5, colframe=red!70!black, boxrule=0.8pt,
  code={\refstepcounter{definition}},
  title={定義 \thedefinition\quad #1}
}
\newtcolorbox{formulabox}[1]{
  mathbox, colback=blue!5, colframe=blue!70!black, boxrule=0.8pt,
  code={\refstepcounter{formula}},
  title={公式 \theformula\quad #1}
}
\newtcolorbox{theorembox}[1]{
  mathbox, colback=blue!5, colframe=blue!70!black, boxrule=1.2pt,
  code={\refstepcounter{theorem}},
  title={定理 \thetheorem\quad #1}
}


\pagestyle{empty}

% ========================================
% ヘッダ（1ページ目の先頭）
% 使い方: \lectureheader{科目}{章番号}{章の名前}{節番号}{節の名前}{項番号}{項の名前}{本時の目標}
% ========================================
\newcommand{\lectureheader}[8]{%
  \setcounter{mathchapter}{#2}%
  \setcounter{section}{#4}%
  \setcounter{subsection}{#6}%
  \noindent\hfill
  年\hspace{2zw}組\hspace{2zw}番\quad 氏名\ \underline{\hspace{8zw}}\par
  \vspace{0.3em}
  \noindent #1\quad 第#2章\quad #3\par
  \noindent{\Large\bfseries 第#4節\quad #5}\hspace{1.5zw}{\large\bfseries #6\quad #7}\par
  \vspace{-0.3em}
  \noindent\rule{\linewidth}{0.8pt}\par
  \vspace{0.3em}
  \begin{itembox}[l]{\textbf{本時の目標}}
    \centering #8
  \end{itembox}
  \vspace{0.3em}
}

% 右のページ（2ページ目）へ移る
\newcommand{\nextpage}{\newpage}

% 見出し: \heading{見出し}
\newcommand{\heading}[1]{%
  \par\noindent{\large\bfseries #1}\par\vspace{0.3em}%
}

% 本時のまとめ
\newenvironment{summarybox}{%
  \begin{itembox}[l]{\textbf{本時のまとめ}}
}{\end{itembox}}

% 左右配置（左：説明・計算，右：図）
% 使い方: \sidebyside[左の幅の割合]{左}{右}
\newcommand{\sidebyside}[3][0.5]{%
  \par\noindent
  \begin{minipage}[t]{#1\linewidth}\vspace{0pt}#2\end{minipage}\hfill
  \begin{minipage}[t]{\dimexpr0.98\linewidth-#1\linewidth\relax}%
    \vspace{0pt}\centering #3\end{minipage}\par
}

% 次回予告（2ページ目の下端）
\newcommand{\nexttime}[1]{%
  \par\noindent\parbox{\linewidth}{%
    \rule{\linewidth}{0.4pt}\par
    次回：#1}\par
}

% ※ B4 に縮小（約87%）すると，17pt は実寸で約14.7pt になる
% ※ jsarticle は \mag で拡大するので，紙面の寸法は truemm で指定する．
%    TikZ の cm も拡大される（1cm → 実寸 約1.7cm）点に注意
% ========================================
\usepackage[dvipdfm, a4paper, margin=15truemm]{geometry}

% ========================================
% 共通プリアンブル（講義MODE・演習MODE）
% パッケージ／解答切り替え／番号／空欄／定義・公式・定理の囲み
% ========================================
\usepackage{amsmath, amssymb, mathtools, bm}
\usepackage{ascmac}
\usepackage[dvipdfmx]{graphicx}
\usepackage[dvipdfmx]{xcolor}
\usepackage{tikz}
\usetikzlibrary{calc, arrows.meta}
\usepackage[most]{tcolorbox}
\usepackage{enumitem}

\linespread{1.25}
\everymath{\displaystyle}
\setlength{\parindent}{0pt}
\setlist[enumerate,1]{label=（\arabic*）, leftmargin=*, labelsep=0.5em}

% ========================================
% 解答表示の切り替え（生徒用：\showanswerfalse ／ 教師用：\showanswertrue）
% ========================================
\newif\ifshowanswer
\showanswertrue

% ========================================
% 番号（「章.節.番号」形式，節ごとにリセット）
% ========================================
% jsarticle には章がないので，章番号は専用カウンタで持つ
\newcounter{mathchapter}
\newcounter{definition}[section]
\newcounter{formula}[section]
\newcounter{theorem}[section]
\renewcommand{\thedefinition}{\arabic{mathchapter}.\arabic{section}.\arabic{definition}}
\renewcommand{\theformula}{\arabic{mathchapter}.\arabic{section}.\arabic{formula}}
\renewcommand{\thetheorem}{\arabic{mathchapter}.\arabic{section}.\arabic{theorem}}

% ========================================
% 書き込み用コマンド
% ========================================
% 空欄（答えの幅＋左右の余白で下線を引く）
% 使い方: \fitblankbf[]{答え}  ※文字列は \text{...} で囲む
\newsavebox{\fitblanksavebox}
\newcommand{\fitblankbf}[2][]{%
  \sbox{\fitblanksavebox}{\ensuremath{\bm{#2}}}%
  \underline{\makebox[\dimexpr\wd\fitblanksavebox+3zw\relax][c]{%
    \ifshowanswer\textcolor{blue}{\usebox{\fitblanksavebox}}%
    \else\phantom{\usebox{\fitblanksavebox}}\fi}}%
}

% 数式中の計算欄（下線なし，答えの幅＋左右2em を確保）
% 使い方: $= \answermath{3 - 1 = 2}$
\newcommand{\answermath}[1]{%
  \ifshowanswer{\color{blue}#1}%
  \else\hspace{2em}\phantom{#1}\hspace{2em}\fi
}

% 記述・計算スペース（薄い枠で高さを確保）
% 使い方: \writespace[高さ]{解答例}
\newcommand{\writespace}[2][5em]{%
  \begin{tcolorbox}[enhanced, colback=white, colframe=black!35,
      boxrule=0.4pt, sharp corners, height=#1, valign=center,
      left=8pt, right=8pt]
    \ifshowanswer\textcolor{blue}{#2}\fi
  \end{tcolorbox}%
}

% ========================================
% 定義（赤系）・公式（青系）・定理（青系・太枠）
% 使い方: \begin{definitionbox}{用語} ... \end{definitionbox}
% ========================================
\tcbset{
  mathbox/.style={
    enhanced, sharp corners,
    left=10pt, right=10pt, top=8pt, bottom=8pt,
    fonttitle=\bfseries,
  }
}
\newtcolorbox{definitionbox}[1]{
  mathbox, colback=red!5, colframe=red!70!black, boxrule=0.8pt,
  code={\refstepcounter{definition}},
  title={定義 \thedefinition\quad #1}
}
\newtcolorbox{formulabox}[1]{
  mathbox, colback=blue!5, colframe=blue!70!black, boxrule=0.8pt,
  code={\refstepcounter{formula}},
  title={公式 \theformula\quad #1}
}
\newtcolorbox{theorembox}[1]{
  mathbox, colback=blue!5, colframe=blue!70!black, boxrule=1.2pt,
  code={\refstepcounter{theorem}},
  title={定理 \thetheorem\quad #1}
}


\pagestyle{empty}

% ========================================
% ヘッダ（1ページ目の先頭）
% 使い方: \lectureheader{科目}{章番号}{章の名前}{節番号}{節の名前}{項番号}{項の名前}{本時の目標}
% ========================================
\newcommand{\lectureheader}[8]{%
  \setcounter{mathchapter}{#2}%
  \setcounter{section}{#4}%
  \setcounter{subsection}{#6}%
  \noindent\hfill
  年\hspace{2zw}組\hspace{2zw}番\quad 氏名\ \underline{\hspace{8zw}}\par
  \vspace{0.3em}
  \noindent #1\quad 第#2章\quad #3\par
  \noindent{\Large\bfseries 第#4節\quad #5}\hspace{1.5zw}{\large\bfseries #6\quad #7}\par
  \vspace{-0.3em}
  \noindent\rule{\linewidth}{0.8pt}\par
  \vspace{0.3em}
  \begin{itembox}[l]{\textbf{本時の目標}}
    \centering #8
  \end{itembox}
  \vspace{0.3em}
}

% 右のページ（2ページ目）へ移る
\newcommand{\nextpage}{\newpage}

% 見出し: \heading{見出し}
\newcommand{\heading}[1]{%
  \par\noindent{\large\bfseries #1}\par\vspace{0.3em}%
}

% 本時のまとめ
\newenvironment{summarybox}{%
  \begin{itembox}[l]{\textbf{本時のまとめ}}
}{\end{itembox}}

% 左右配置（左：説明・計算，右：図）
% 使い方: \sidebyside[左の幅の割合]{左}{右}
\newcommand{\sidebyside}[3][0.5]{%
  \par\noindent
  \begin{minipage}[t]{#1\linewidth}\vspace{0pt}#2\end{minipage}\hfill
  \begin{minipage}[t]{\dimexpr0.98\linewidth-#1\linewidth\relax}%
    \vspace{0pt}\centering #3\end{minipage}\par
}

% 次回予告（2ページ目の下端）
\newcommand{\nexttime}[1]{%
  \par\noindent\parbox{\linewidth}{%
    \rule{\linewidth}{0.4pt}\par
    次回：#1}\par
}


\begin{document}

\lectureheader{数学Ⅱ}{6}{微分法と積分法}{1}{微分係数と導関数}{1}{微分係数}
  {平均変化率の代数的なイメージを持とう}

% 同じ節の2枚目以降は，前回までの番号を引き継ぐ
% \setcounter{definition}{1}

% ============ 1ページ目（表・左） ============
\begin{reviewbox}{中学数学：1次関数の変化の割合}
  1次関数 $y = ax + b$ では,
  \[
    \text{変化の割合}
    = \frac{\fitblankbox{y \text{ の増加量}}}{\fitblankbox{x \text{ の増加量}}}
    = \fitblankbf[]{a}
  \]
  であり, 変化の割合は常に一定である.

  \vspace{0.5em}
  例）$y = 2x + 1$ で, $x$ が $1$ から $4$ まで増加するとき\par
  \hspace{1zw}$x$ の増加量 $= \answermath{3}$,\quad
  $y$ の増加量 $= \answermath{6}$\quad より\quad
  変化の割合 $= \answermath{\frac{6}{3} = 2}$

  \vspace{0.3em}
  では, グラフが直線でない $y = x^2$ などでは, どう考える?
\end{reviewbox}

\topic{A}{平均変化率}

\heading{$x$ の変化量と $y$ の変化量}
関数 $y = f(x)$ で, $x$ の値を $a$ から $b$ まで変化させる.

\vspace{0.3em}
\sidebyside[0.44]{%
  $x$ の変化量\par
  \hspace{1zw}$=$ \fitblankbf[]{b - a}

  \vspace{0.8em}
  $y$ の変化量\par
  \hspace{1zw}$=$ \fitblankbf[]{f(b) - f(a)}

  \vspace{0.8em}
  変化量は, いつも\par
  （\fitblankbf[]{\text{後}}）$-$（\fitblankbf[]{\text{前}}）
}{%
  \begin{tikzpicture}[x=0.95cm, y=0.95cm, >=Stealth,
      font=\footnotesize]
    \draw[->] (-0.2,0) -- (4.0,0) node[right] {$x$};
    \draw[->] (0,-0.2) -- (0,4.0) node[above] {$y$};
    \node[anchor=north east] at (0,0) {$O$};
    % 直線 AB
    \draw[gray] (0.3,0.62) -- (3.4,3.33);
    % 曲線 y = f(x)
    \draw[thick, domain=0.2:3.4, samples=60]
      plot (\x, {0.25*\x*\x + 0.6}) node[above] {$y = f(x)$};
    % 補助線
    \draw[dotted] (0.8,0) -- (0.8,0.76) -- (0,0.76);
    \draw[dotted] (2.8,0) -- (2.8,2.56) -- (0,2.56);
    \node[below] at (0.8,0) {$a$};
    \node[below] at (2.8,0) {$b$};
    \node[left] at (0,0.76) {$f(a)$};
    \node[left] at (0,2.56) {$f(b)$};
    % 変化量
    \draw[<->, semithick] (0.8,0.76) -- (2.8,0.76)
      node[midway, below] {$x$ の変化量};
    \draw[<->, semithick] (2.8,0.76) -- (2.8,2.56)
      node[midway, right, align=left] {$y$ の\\変化量};
    % 点
    \fill (0.8,0.76) circle (1.5pt) node[above left] {A};
    \fill (2.8,2.56) circle (1.5pt) node[above left] {B};
  \end{tikzpicture}
}

\vfill
\nextpage

% ============ 2ページ目（表・右） ============\begin{definitionbox}{平均変化率}
  関数 $y = f(x)$ において, $x$ の値が $a$ から $b$ まで変化するとき,
  \[
    \frac{\fitblankbox{y \text{ の変化量}}}{\fitblankbox{x \text{ の変化量}}}
    = \frac{f(b) - f(a)}{b - a}
  \]
  を, $x = a$ から $x = b$ までの $f(x)$ の \textbf{平均変化率} という.
\end{definitionbox}

\vspace{0.5em}
グラフで見ると, 平均変化率は\par
2点 A, B を結ぶ直線の \fitblankbf[]{\text{傾き}} である.

\vfill
\heading{具体例で確かめる}
関数 $f(x) = x^2$ について, $x = 1$ から $x = 3$ までの平均変化率を求める.

\vspace{0.3em}
\sidebyside[0.56]{%
  $x$ の変化量\par
  \hspace{1zw}$= \answermath{3 - 1 = 2}$

  \vspace{0.8em}
  $y$ の変化量\par
  \hspace{1zw}$= \answermath{f(3) - f(1) = 9 - 1 = 8}$

  \vspace{0.8em}
  平均変化率\par
  \hspace{1zw}$= \answermath{\frac{8}{2} = 4}$
}{%
  \begin{tikzpicture}[x=0.9cm, y=0.32cm, >=Stealth, font=\footnotesize]
    \draw[->] (-0.4,0) -- (3.8,0) node[right] {$x$};
    \draw[->] (0,-0.8) -- (0,10.8) node[above] {$y$};
    \node[anchor=north east] at (0,0) {$O$};
    \draw[gray] (0.55,-0.8) -- (3.25,10);
    \draw[thick, domain=-0.3:3.2, samples=60]
      plot (\x, {\x*\x}) node[above] {$y = x^2$};
    \draw[dotted] (1,0) -- (1,1) -- (0,1);
    \draw[dotted] (3,0) -- (3,9) -- (0,9);
    \node[below] at (1,0) {$1$};
    \node[below] at (3,0) {$3$};
    \node[left] at (0,1) {$1$};
    \node[left] at (0,9) {$9$};
    \draw[<->, semithick] (1,1) -- (3,1);
    \draw[<->, semithick] (3,1) -- (3,9);
    \ifshowanswer
      \node[above] at (2,1) {\answerstyle{$2$}};
      \node[right] at (3,5) {\answerstyle{$8$}};
    \fi
    \fill (1,1) circle (1.5pt);
    \fill (3,9) circle (1.5pt);
  \end{tikzpicture}
}

\vfill
平均変化率が $4$ であるとは, どういうことか.\par
言葉で書いてみよう.
\writespace[4em]{$x$ が $1$ 増えるあたり, $y$ は平均して $4$ 増える.}

\vfill
\nextpage

% ============ 3ページ目（裏・左） ============
\begin{summarybox}
  \begin{enumerate}[label=\textbf{\arabic*}, leftmargin=1.5zw, itemsep=0.5em]
    \item \textbf{用語}\quad 「平均変化率」とは何か.\par
      「変化量」という言葉を使って説明しよう.
      \writespace[4.5em]{$x$ の値が $a$ から $b$ まで変化するときの,
        $y$ の変化量を $x$ の変化量で割った値.}

    \item \textbf{式}\quad $x = a$ から $x = b$ までの $f(x)$ の平均変化率を式で書こう.
      \writespace[4em]{\centering $\dfrac{f(b) - f(a)}{b - a}$\par}
      分子 $f(b) - f(a)$ は \fitblankbf[]{y \text{ の変化量}} を,\par
      分母 $b - a$ は \fitblankbf[]{x \text{ の変化量}} を表す.

    \item \textbf{グラフ}\quad 下の図に, 2点 A$(a,\,f(a))$, B$(b,\,f(b))$ と,
      $x$ の変化量, $y$ の変化量をかき込もう.

      \vspace{0.3em}
      \begin{center}
      \begin{tikzpicture}[x=0.95cm, y=0.8cm, >=Stealth, font=\footnotesize]
        \draw[->] (-0.2,0) -- (4.2,0) node[right] {$x$};
        \draw[->] (0,-0.2) -- (0,3.8) node[above] {$y$};
        \node[anchor=north east] at (0,0) {$O$};
        \draw[thick, domain=0.2:3.4, samples=60]
          plot (\x, {0.25*\x*\x + 0.6}) node[above] {$y = f(x)$};
        \ifshowanswer
          \begin{scope}[blue, semithick]
            \draw (0.3,0.62) -- (3.4,3.33);
            \draw[dotted] (0.8,0) -- (0.8,0.76) -- (0,0.76);
            \draw[dotted] (2.8,0) -- (2.8,2.56) -- (0,2.56);
            \node[below] at (0.8,0) {$a$};
            \node[below] at (2.8,0) {$b$};
            \node[left] at (0,0.76) {$f(a)$};
            \node[left] at (0,2.56) {$f(b)$};
            \draw[<->] (0.8,0.76) -- (2.8,0.76)
              node[midway, below] {\answerstyle{$x$ の変化量}};
            \draw[<->] (2.8,0.76) -- (2.8,2.56)
              node[midway, right, align=left] {\answerstyle{$y$ の}\\\answerstyle{変化量}};
            \fill (0.8,0.76) circle (1.5pt) node[above left] {A};
            \fill (2.8,2.56) circle (1.5pt) node[above left] {B};
          \end{scope}
        \fi
      \end{tikzpicture}
      \end{center}
      平均変化率は, 直線 AB の \fitblankbf[]{\text{傾き}} を表す.

    \item \textbf{中学との違い}\par
      1次関数では, 平均変化率（変化の割合）は $x$ の区間によらず
      \fitblankbf[]{\text{一定}} である.\par
      一方, $f(x) = x^2$ のような関数では, 平均変化率は\par
      $x$ の区間によって \fitblankbf[]{\text{変わる}}.
  \end{enumerate}
\end{summarybox}

\vfill
\nextpage

% ============ 4ページ目（裏・右） ============
\heading{確認問題}
次の関数について, 指定された区間の平均変化率を求めよ.

\begin{enumerate}[label=\textbf{問\arabic*}, leftmargin=3zw, itemsep=0.8em]
  \item $f(x) = x^2$\quad （$x = 2$ から $x = 5$ まで）
    \writespace[4.5em]{$\dfrac{f(5) - f(2)}{5 - 2} = \dfrac{25 - 4}{3} = 7$}

  \item $f(x) = x^2$\quad （$x = -3$ から $x = -1$ まで）
    \writespace[4.5em]{$\dfrac{f(-1) - f(-3)}{-1 - (-3)} = \dfrac{1 - 9}{2} = -4$}

  \item $f(x) = x^2 - 3x$\quad （$x = 1$ から $x = 4$ まで）
    \writespace[4.5em]{$f(4) = 4$, $f(1) = -2$ より\quad
      $\dfrac{4 - (-2)}{4 - 1} = \dfrac{6}{3} = 2$}

  \item $f(x) = 3x + 1$ について, 次の区間の平均変化率を求め,\par
    気づいたことを書け.\par
    （1）$x = 1$ から $x = 4$ まで\quad （2）$x = -2$ から $x = 0$ まで
    \writespace[7em]{（1）$\dfrac{13 - 4}{3} = 3$\quad
      （2）$\dfrac{1 - (-5)}{2} = 3$\par
      区間によらず $3$（直線の傾き）で一定である.}
\end{enumerate}

\vfill
\nexttime{$b$ を $a$ に限りなく近づけると, 平均変化率はどうなるだろうか.}

\end{document}
